\documentclass[12pt, a4paper]{article}

% --- 常用宏包 ---
\usepackage[UTF8]{ctex}
\usepackage{amsmath, amsthm, amssymb, amsfonts}
\usepackage{geometry}
\usepackage{fancyhdr}
\usepackage{enumerate}
\usepackage{graphicx}
\usepackage{hyperref}
\usepackage{bm}
\usepackage{tcolorbox}

% --- 页面设置 ---
\geometry{left=2.5cm, right=2.5cm, top=2.5cm, bottom=2.5cm}
\pagestyle{fancy}
\fancyhf{}
\rhead{\today}
\lhead{近世代数作业}
\cfoot{\thepage}

% --- 环境定义 ---
\newtcolorbox{problem}[1]{
    colback=gray!5!white,
    colframe=gray!75!black,
    fonttitle=\bfseries,
    title=题目 #1
}

\newenvironment{solution}{
    \begin{proof}[\textbf{解}]
}{
    \end{proof}
}

% --- 个人信息 ---
\title{近世代数作业 02}
\author{李睿阳 (524120910059)}
\date{\today}

\begin{document}

\maketitle

\begin{problem}{1}
    求加法循环群 $\mathbb{Z}_{10}$ 的所有生成元和所有子群。
\end{problem}
\begin{solution}
    $\mathbb{Z}_{10}$ 的生成元 $x$ 满足 $\gcd(10,x)=1$，生成元为 $1, 3, 7, 9$。

    $\mathbb{Z}_{10}$ 的所有子群为：
    \begin{itemize}
        \item 10 阶子群：$\langle 1 \rangle = \{0, 1, 2, 3, 4, 5, 6, 7, 8, 9\}$
        \item 5 阶子群：$\langle 2 \rangle = \{0, 2, 4, 6, 8\}$
        \item 2 阶子群：$\langle 5 \rangle = \{0, 5\}$
        \item 1 阶子群：$\langle 0 \rangle = \{0\}$
    \end{itemize}
\end{solution}

\begin{problem}{2}
    找出 $\mathbb{Z}_{11}^*$ 的一个生成元，并将所有元素表示为生成元的幂。
\end{problem}

\begin{solution}
    $\mathbb{Z}_{11}^* = \{1, 2, 3, 4, 5, 6, 7, 8, 9, 10\}$。

    选取生成元 $2$，各个元素的阶分别为：
    \begin{itemize}
        \item $2^1 \equiv 2 \pmod{11}$
        \item $2^2 \equiv 4 \pmod{11}$
        \item $2^3 \equiv 8 \pmod{11}$
        \item $2^4 \equiv 16 \equiv 5 \pmod{11}$
        \item $2^5 \equiv 10 \pmod{11}$
        \item $2^6 \equiv 20 \equiv 9 \pmod{11}$
        \item $2^7 \equiv 18 \equiv 7 \pmod{11}$
        \item $2^8 \equiv 14 \equiv 3 \pmod{11}$
        \item $2^9 \equiv 6 \pmod{11}$
        \item $2^{10} \equiv 1 \pmod{11}$
    \end{itemize}
\end{solution}

\begin{problem}{3}
    求商群 $\mathbb{Z}_{12}/\langle 4 \rangle$ 的各个元素及加法运算表。
\end{problem}

\begin{solution}
    $\langle 4 \rangle = \{0, 4, 8\}$ 为 $\mathbb{Z}_{12}$ 的子群，商群 $\mathbb{Z}_{12} / \langle 4 \rangle$ 的元素为：
    \begin{itemize}
        \item $0 + \langle 4 \rangle = \{0, 4, 8\}$
        \item $1 + \langle 4 \rangle = \{1, 5, 9\}$
        \item $2 + \langle 4 \rangle = \{2, 6, 10\}$
        \item $3 + \langle 4 \rangle = \{3, 7, 11\}$
    \end{itemize}

    加法运算表：
    \begin{table}[h]
        \centering
        \begin{tabular}{c|cccc}
            + & $0 + \langle 4 \rangle$ & $1 + \langle 4 \rangle$ & $2 + \langle 4 \rangle$ & $3 + \langle 4 \rangle$ \\ \hline
            $0 + \langle 4 \rangle$ & $0 + \langle 4 \rangle$ & $1 + \langle 4 \rangle$ & $2 + \langle 4 \rangle$ & $3 + \langle 4 \rangle$ \\
            $1 + \langle 4 \rangle$ & $1 + \langle 4 \rangle$ & $2 + \langle 4 \rangle$ & $3 + \langle 4 \rangle$ & $0 + \langle 4 \rangle$ \\
            $2 + \langle 4 \rangle$ & $2 + \langle 4 \rangle$ & $3 + \langle 4 \rangle$ & $0 + \langle 4 \rangle$ & $1 + \langle 4 \rangle$ \\
            $3 + \langle 4 \rangle$ & $3 + \langle 4 \rangle$ & $0 + \langle 4 \rangle$ & $1 + \langle 4 \rangle$ & $2 + \langle 4 \rangle$
        \end{tabular}
    \end{table}
\end{solution}
    
\begin{problem}{4}
    证明：群的两个正规子群的交是正规子群。
\end{problem}

\begin{solution}
    易得，两个子群的交仍为子群。

    由正规子群的定义： $H \triangleleft G$ 要求 $\forall g \in G, \forall x \in H$，有 $gxg^{-1} \in H$。

    设 $H, K$ 为群 $G$ 的正规子群。对于 $\forall x \in H \cap K$，由此可知 $x \in H$ 且 $x \in K$。
    
    因为 $H, K$ 均是 $G$ 的正规子群，对于 $\forall g \in G$，有：
    \begin{itemize}
        \item 由于 $x \in H$ 且 $H \triangleleft G$，所以 $gxg^{-1} \in H$；
        \item 由于 $x \in K$ 且 $K \triangleleft G$，所以 $gxg^{-1} \in K$。
    \end{itemize}
    
    由 $gxg^{-1} \in H$ 与 $gxg^{-1} \in K$ 得到 $gxg^{-1} \in H \cap K$。
    
    由正规子群的定义，$H \cap K$ 也是 $G$ 的正规子群。
\end{solution}

\begin{problem}{5}
    设 $f: \mathbb{Z}_{12} \to \mathbb{Z}_3$ 为群同态，且 $f(1) = 2$。
    \begin{enumerate}[(1)]
        \item 证明：$f$ 是满同态；
        \item 求 $\ker(f)$；
        \item 列出 $\mathbb{Z}_{12}/\ker(f)$ 的所有陪集；
        \item 给出由 $\mathbb{Z}_{12}/\ker(f) \to \mathbb{Z}_3$ 的同构。
    \end{enumerate}
\end{problem}

\begin{solution}
    \begin{enumerate}[(1)]
        \item 由 $f$ 是群同态，对于 $\forall n \in \mathbb{Z}_{12}$，有 $f(n) = f(n \cdot 1) = n \cdot f(1) = 2n \pmod{3}$。
        可得到：
        \begin{itemize}
            \item $f(0) = 0$
            \item $f(1) = 2$
            \item $f(2) = 2 \times 2 \equiv 1 \pmod{3}$
        \end{itemize}
        由于 $\text{Im}(f) = \{0, 1, 2\} = \mathbb{Z}_3$， $f$ 是满同态。

        \item $\ker(f) = \{x \in \mathbb{Z}_{12} \mid f(x) = 0\} = \{x \in \mathbb{Z}_{12} \mid 2x \equiv 0 \pmod{3}\}$。
        
            则 $\ker(f) = \{0, 3, 6, 9\}$。

        \item 商群 $\mathbb{Z}_{12}/\ker(f)$ 的陪集为：
        \begin{itemize}
            \item $0 + \ker(f) = \{0, 3, 6, 9\}$
            \item $1 + \ker(f) = \{1, 4, 7, 10\}$
            \item $2 + \ker(f) = \{2, 5, 8, 11\}$
        \end{itemize}

        \item
        同构为；
        \begin{itemize}
            \item $\bar{f}(\{0, 3, 6, 9\}) = f(0) = 0$
            \item $\bar{f}(\{1, 4, 7, 10\}) = f(1) = 2$
            \item $\bar{f}(\{2, 5, 8, 11\}) = f(2) = 1$
        \end{itemize}
    \end{enumerate}
\end{solution}

\end{document}
